\chapter{Names Across Volume Modes}

Path names come from directory record keys, stored as UTF-8 byte
strings. The chapter develops one positive finding (verbatim
preservation) and one deliberate scope limit (lookup-by-name is
not claimed).

\section{Verbatim Preservation}

\textbf{Observation.} The Rust body decoder emits the bytes the
directory key stores. The parser does not normalise, case-fold, or
interpret. POSIX traversal of the same volume returns the same
bytes. The two agree byte for byte on both case-insensitive and
case-sensitive APFS volumes.

Anything more elaborate --- ``emit the canonical form,'' ``convert
to NFC'' --- would be wrong on some volume and would silently
rename files in the output.

\section{Two Volume Modes Exist}

APFS volumes come in two case-and-normalisation modes that the
indexer must handle.

\begin{description}
  \item[\texttt{APFS} (case-insensitive).]
    \texttt{case\_insensitive = true},
    \texttt{normalization\_insensitive = false}. The volume's
    incompatible-feature mask sets the case-insensitive bit only.
    The volume is also normalisation-insensitive for lookup, but
    APFS reserves the explicit incompat bit for the case-sensitive
    variant. The asymmetry is a quirk of the on-disk format, not a
    statement about behaviour.
  \item[\texttt{Case-sensitive APFS}.]
    \texttt{case\_insensitive = false},
    \texttt{normalization\_insensitive = true}.
\end{description}

Both modes appear on real systems, and both have to be supported.

\section{What The Modes Affect}

The modes affect \emph{lookup}: when the kernel resolves a path it
consults the volume's case and normalisation behaviour to decide
whether two names refer to the same file. The modes do not affect
row \emph{enumeration}, which walks directory keys and returns
their bytes. The parser emits correct rows on both modes without
performing the lookup-equivalence test, because lookup is the
running system's responsibility.

\subsection{The Lookup Oracle}

\textbf{Observation.} Across the two volume modes, APFS lookup
behaves as follows. On the EX-20 same-run fixture using
\texttt{os.O\_CREAT | os.O\_EXCL}:

\begin{center}
\begin{tabular}{lll}
\toprule
operation & CI volume & CS volume \\
\midrule
create \texttt{plain.txt}                          & accept & accept \\
create \texttt{CaseName.txt}                       & accept & accept \\
\texttt{casename.txt} after \texttt{CaseName.txt}  & EEXIST & accept \\
create NFC \texttt{caf\'e.txt}                     & accept & accept \\
NFD \texttt{caf\'e.txt} after NFC sibling          & EEXIST & EEXIST \\
\bottomrule
\end{tabular}
\end{center}

The case-insensitive volume rejects the lowercase form because lookup
considers the two names equivalent. Both volumes reject the
decomposed Unicode form after the precomposed sibling because both
treat NFC and NFD as the same name for lookup. This is what APFS does;
no APFS-specific knowledge in the indexer is needed to produce the
result, because the kernel did the work.

\subsection{Rust vs POSIX Equality}

\textbf{Observation.} The Rust-emitted path bytes match
POSIX-traversal path bytes exactly on both volumes:
\texttt{path.encode("utf-8").hex()} produces the same string on
both sides for every entry. CI corpus has four entries; CS has
five (the case-distinct sibling also appears); equality holds for
every row. Both sides walk the directory key storage; POSIX walks
through the kernel, Rust walks the FS tree directly; neither
applies normalisation, so the bytes are the same bytes.

\section{Lookup-by-Name Is Not Claimed}

The APFS lookup primitive is the directory hash: CRC32C over the
name's UTF-32LE codepoints, with NFD normalisation applied first,
and Unicode case-fold applied before normalisation on
case-insensitive volumes. Two names with the same hash collide
within a directory; APFS resolves the collision by storing both
keys and walking them.

A \emph{find this exact path} feature must reproduce that hash. The
project has not validated a hash implementation against a same-run
oracle, and until it does the parser refuses to find files by name.
Row enumeration powers the treemap, sort-by-size, drill-down, and
context-menu operations on rows the user already sees; a silently
wrong lookup-by-name returns the wrong file for the user's query.
The cost asymmetry dictates the scope asymmetry.

The substring search the native shell ships (chapter~13) operates
on already-enumerated names: it does not perform an APFS-hash
lookup and therefore inherits no risk from the missing hash.

\section{Names and Aggregates}

Directory aggregates (chapter~8) key on inode identity, not on
path. Two paths that hard-link to the same inode contribute that
inode's size once to a directory total, regardless of whether their
names are casewise or normalisationally distinct. Hard-link
unification reads SIBLING\_LINK and SIBLING\_MAP to find every name
pointing at a given inode; names are emitted verbatim in each row,
the inode is recorded as the row's \texttt{file\_id}, and the
aggregate keys on the file\_id.

\section{What A Future Lookup-by-Name Would Need}

\begin{enumerate}
  \item A same-run hash oracle: take a constructed name, run
    APFS's hash through a controlled API (the kernel's
    \texttt{O\_EXCL} accept/reject is the cheapest indicator),
    produce positive and negative cases for the hash equivalence
    class.
  \item NFD normalisation tables and Unicode case-fold tables
    matching the macOS version the volume was written on
    (tables change across Unicode revisions).
  \item Collision handling: a directory may store multiple keys
    with the same hash; the lookup walks all of them.
  \item Volume-mode dispatch: case-insensitive applies case-fold
    before normalisation; case-sensitive does not.
\end{enumerate}
